\documentclass[12pt]{article}
\usepackage[T1]{fontenc}
\usepackage[utf8]{inputenc}
\usepackage{lmodern}
\usepackage{textcomp}
\usepackage{enumitem}
\usepackage{geometry}
\geometry{margin=1in}
\begin{document}
\begin{center}
Answer Key - Sociology - Graduate Level Assessment 25 \
\small (Answers are ordered exactly as in the question paper.)
\end{center}

\section*{Multiple Choice}
\begin{enumerate}[label=\arabic*.]
\item \textbf{Answer:} C
\\ \textit{Options:} A) genetics; B) evolution; C) height; D) brain size
\\ \textit{Note:} Height was not considered a significant factor in the naturalistic explanations of race during the nineteenth century, as the focus was primarily on other physical traits such as skull shape and skin color, which were erroneously linked to intellectual and moral capacities.
\item \textbf{Answer:} C
\\ \textit{Options:} A) lacking exclusive use of a bath or shower; B) living in housing with serious structural defects; C) buying fewer than twenty DVDs in the previous year; D) going without a week's holiday in the previous year
\\ \textit{Note:} The correct answer is "buying fewer than twenty DVDs in the previous year" because Townsend's indicators of relative deprivation focus on essential needs and social participation rather than specific consumer behaviors or the quantity of non-essential items purchased.
\item \textbf{Answer:} B
\\ \textit{Options:} A) charismatic authority; B) rational-legal authority; C) traditional authority; D) value-rational authority
\\ \textit{Note:} Weber's concept of rational-legal authority posits that the legitimacy of the state's monopoly on the use of force arises from a system of rules and laws that are accepted by the populace, distinguishing it from other forms of authority based on tradition or charisma.
\item \textbf{Answer:} A
\\ \textit{Options:} A) individuals buying welfare privately, with some means-tested benefits; B) regular benefit payments to men as earners of the 'family wage'; C) a universalist system of welfare for all, regardless of income; D) decommodifying social welfare through state provision
\\ \textit{Note:} The market model of state welfare emphasizes the role of individual responsibility in securing welfare through private purchases and aims to supplement this with targeted means-tested benefits for those who cannot afford to do so, reflecting a belief in market efficiency and personal choice in welfare provision.
\item \textbf{Answer:} D
\\ \textit{Options:} A) sexual characteristics are the biological determinants of gender; B) heterosexuality and homosexuality are essential, opposing identities; C) the 'two-sex' model replaced the 'one-sex' model in the eighteenth century; D) gender is performed through bodily gestures and styles to create 'sex'
\\ \textit{Note:} Judith Butler's theory posits that gender is not an innate quality but rather a social construct that is enacted and reinforced through repeated bodily gestures and performances, thereby shaping our understanding of 'sex' as a category.
\end{enumerate}

\section*{True / False}
\begin{enumerate}[label=\arabic*.]
\setcounter{enumi}{5}
\item \textbf{Answer:} False
\\ \textit{Options:} A) True; B) False
\\ \textit{Note:} Thatcher's government implemented policies that prioritized reducing public spending and welfare dependency, which often resulted in reduced financial support for single parents, students, and the unemployed.
\item \textbf{Answer:} True
\\ \textit{Options:} A) True; B) False
\\ \textit{Note:} Chodorow (1978) posits that while both genders initially develop a strong emotional bond with their mothers, cultural and societal expectations lead boys to disengage from this maternal attachment as they mature, resulting in a differentiated emotional development between boys and girls.
\item \textbf{Answer:} False
\\ \textit{Options:} A) True; B) False
\\ \textit{Note:} Researchers have access to a variety of official statistics on domestic violence from government agencies and organizations, which enables them to study the issue effectively.
\item \textbf{Answer:} False
\\ \textit{Options:} A) True; B) False
\\ \textit{Note:} Decarceration primarily advocates for reducing reliance on imprisonment and promoting alternative forms of punishment and rehabilitation, rather than increasing its use.
\item \textbf{Answer:} False
\\ \textit{Options:} A) True; B) False
\\ \textit{Note:} The human relations approach emphasizes the importance of social factors, employee motivation, and interpersonal relationships over strict control and discipline in fostering high productivity.
\end{enumerate}

\section*{Long Form Response}
\begin{enumerate}[label=\arabic*.]
\setcounter{enumi}{10}
\item \textbf{Answer:} In sociology, the concept of community encompasses various forms, including geographic, interest-based, and virtual communities, all of which foster social connections and shared identities among their members. Sociological communities formed by unpopular lecturers may arise from a shared experience of marginalization within the academic environment; however, they often lack the cohesive features that characterize more traditional communities, such as mutual support and collective purpose. This lack of cohesion can stem from the negative perceptions associated with the lecturers, which may inhibit the development of a strong communal identity. Consequently, while these groups may function as informal networks, they do not fulfill the criteria of a new form of community, as they do not cultivate the same level of engagement or shared values typically observed in established communities. Thus, the sociological analysis reveals that the dynamics of such communities are more reflective of social isolation than of a robust communal bond.
\\ \textit{Note:} correct because it accurately identifies the various forms of community in sociology while critically assessing how the lack of cohesive features and shared purpose in communities formed by unpopular lecturers limits their classification as a true community.
\item \textbf{Answer:} Émile Durkheim's concept of social facts refers to the norms, values, structures, and institutions that exist outside the individual but exert a powerful influence on behavior and social interactions. He posited that sociology should focus on these social facts because they represent the collective aspects of social life that shape individual actions and societal cohesion. By analyzing social facts, Durkheim believed sociologists could uncover the underlying forces that maintain social order and contribute to the functioning of society. Their significance lies in understanding how these external factors contribute to social stability, collective consciousness, and the overall structure of society, thereby providing a framework for examining social phenomena beyond individual psychology.
\\ \textit{Note:} correct because it accurately describes Durkheim's view that social facts are external, collective forces shaping individual behavior and social cohesion, emphasizing their importance for understanding the structures that maintain societal order.
\end{enumerate}

\vfill
\noindent\textit{End of Answer Key}
\end{document}